\scanpage{135}
\mspage{V.7}
% archive sheet number: 134
% manuscript lines 1-6
\deleted{\unsure{et soit}}

On peut supposer $X$ affine, et on peut trouver une suite d'\unsure{élts|éléments}
\[
t_1,\ldots,t_r\in\Gamma(X,\unsure{\mathcal O|O}_X)
\]
telle que leurs images dans $k(x)$ \unsure{font|forment} une base de transc. sur $k$. Les $t_i$ définissent un morphisme
\[
h:X\mathbin{\unsure{-|\to}}Y=S[T_1,\ldots,T_r],
\]
% manuscript lines 7-11
et si $y=h(x)$, on voit que $k(x)/k(y)$ est fini, donc que $x$ est fermé dans $h^{-1}(y)$. On a donc
\[
\dim_x h^{-1}(y)=\dim\underline{\mathcal O}_{X_0,x}=n-r,
\]
et quitte à restreindre $X$, on peut donc supposer que les fibres de $X$ sur $Y$ sont de dim $\leq n-r$.

% manuscript lines 12-15
Soit alors $S_1$ le localisé strict de $Y$ en le \deleted{le pt fermé} pt géométrique de $Y$ défini par le pt géom $\bar x$ de $X$, $X_1=X\times_Y S_1$, $U_1=U\times_Y S_1$. Alors $X_1$ est affine

% lower-left diagram: entire diagram cancelled by diagonal strokes
\diagram{\begin{tikzcd}[ampersand replacement=\&,column sep=small,row sep=small]
\&\deleted{X'}\arrow[d]\&\deleted{U'}\arrow[l,no head]\arrow[d]\&\deleted{U'_{\unsure{\bar\eta|\bar y}}}\arrow[d,no head]\&\deleted{\unsure{U'\otimes_{S_1}\bar L}}\arrow[l,no head]\arrow[d,no head]\\
\deleted{X}\arrow[d]\&\deleted{X_1}\arrow[l]\arrow[d]\&\deleted{\unsure{U_1}}\arrow[d,"\deleted{\text{\unsure{pas cart.}}}"]\\
\deleted{S[T_1,\ldots,T_r]=Y}\arrow[d]\&\deleted{S_1}\arrow[l]\arrow[d,"\deleted{\unsure{\varphi}}"]\&\deleted{S_1}\arrow[l,equal]\arrow[d,no head]\&\deleted{S_{1\bar\eta}}\arrow[l,no head]\arrow[d,no head]\&\deleted{\bar L}\arrow[l,no head]\\
\deleted{S}\&\deleted{S}\&\deleted{S}\arrow[l,no head]\&\deleted{\bar\eta}\arrow[l,no head]
\end{tikzcd}}
% lower-right revised diagram
\diagram{\begin{tikzcd}[ampersand replacement=\&,column sep=small,row sep=small]
\&X'\arrow[d]\&U'\arrow[l]\arrow[d,no head]\&U'_{\bar\eta}\arrow[l]\arrow[dd,"\varphi"]\&U'\times_{S_1}\bar\lambda\arrow[l,no head]\arrow[dd]\\
X\arrow[d]\&X_1\arrow[l]\arrow[d]\&U_1\arrow[l]\arrow[d]\\
S[T_1,\ldots,T_r]=Y\arrow[d]\&S_1\arrow[l]\&S_1\arrow[l,equal]\arrow[d]\&S_{1\bar\eta}\deleted{\unsure{=\lambda}}\arrow[l]\arrow[d,no head]\&\deleted{\unsure{\bar L}}\bar\lambda\arrow[l,no head]\\
S\&\deleted{\unsure{\eta}}\arrow[u]\&S\&\deleted{\unsure{\eta}}\bar\eta\arrow[l]\&\deleted{\unsure{\longrightarrow|-}}
\end{tikzcd}}
% manuscript lines 16-20
\unsure{t.f.|l.f.} sur le trait \unsure{st.} hensélien $S_1$, de dim relative $n-r$, $U_1$ en est un ouvert affine, $X'$ est le localisé \inserted{strict} \uline{de} $X_1$ en le pt géom $\bar x$, $U'$ est l'image inverse de $U_1$ dans $X'$. Soit $L$ le corps des fractions de $S_1$ qui est un
